\chapter{Algorithms for IMDb Rating Prediction}\label{chap:algo}

\section{Overview}

This chapter presents the algorithms used in the IMDb rating prediction system. The system combines SBERT (Sentence Transformer) for semantic text representation with XGBoost for regression to predict movie ratings from scripts and metadata.

The algorithms presented here include:
\begin{itemize}
    \item SBERT embedding generation with chunking strategy for long scripts
    \item XGBoost model training with class balancing
    \item Complete prediction pipeline
\end{itemize}

\section{Referring Functions and Algorithms}

Use \texttt{\textbackslash ref} command preceded with \textbf{Function} or \textbf{Algorithm} to refer them. For example, Function~\ref{func:sbert_embedding} and Algorithm~\ref{alg:train_xgboost}. See the LaTeX source code to see how to refer to functions and algorithms.

\section{SBERT Embedding Generation}

Movie scripts can be very long (often exceeding 10,000 words), which is beyond the token limit of SBERT models ($\sim$256 tokens). To handle this, we use a chunking strategy with overlapping windows and average pooling.

Function~\ref{func:sbert_embedding} shows the complete process of generating SBERT embeddings from movie scripts.

%MARGIN FOR ALGORITHM
\IncMargin{1em}
%START OF FUNCTION
\begin{func}
\begin{algorithm}[!p]
%%%%SETTING KEYWORDS FOR STYLING THEM IN THE FUNCTION%%%%
%%%%%%%%%

\SetKwFunction{chunktext}{ChunkText}
\SetKwFunction{getembedding}{GetSBERTEmbedding}

%%%%%%%%%
%Usage of Input,Output, Data( Persistent data like database)
%%%%%%
\KwInput{ Movie script text $T$ }
\KwOutput{ 384-dimensional embedding vector $E$ }
\KwData{ SBERT model $M$, chunk size $C=256$, overlap $O=50$ }
%%%%LINEGAPE
\BlankLine
%%%%FOR LOOP BLOCK
$chunks \gets \chunktext{T, C, O}$\;
$E \gets \text{zero vector of size } 384$\;
\For{ \upshape{each} $chunk$ \upshape{in} $chunks$ }
{
    $e_{chunk} \gets M.encode(chunk)$\;
    $E \gets E + e_{chunk}$\;
}
$E \gets E / |chunks|$\; \tcp*{Average pooling}
\Return $E$
%%%%%CAPTIONING THE FUNCTION
\caption{GenerateSBERTEmbedding$(T)$}
\label{func:sbert_embedding}
%END OF FUNCTION
\end{algorithm}
\end{func}
%MARGIN FOR ALGORITHM
\DecMargin{1em}


The key aspects of this algorithm are:

\begin{itemize}
    \item \textbf{Chunking Strategy}: Scripts are split into 256-word chunks with 50-word overlap to maintain semantic continuity
    \item \textbf{SBERT Encoding}: Each chunk is encoded independently using the \texttt{all-MiniLM-L6-v2} model
    \item \textbf{Average Pooling}: Multiple chunk embeddings are averaged to produce a single document-level representation
    \item \textbf{Fixed Dimension}: Output is always a 384-dimensional vector regardless of script length
\end{itemize}

The overlapping chunks ensure that important context is preserved across chunk boundaries, while average pooling produces a robust document-level embedding that represents the entire script's semantic content.

\section{XGBoost Model Training}

Algorithm~\ref{alg:train_xgboost} presents the complete training pipeline including data splitting, feature engineering, sample weighting for class balancing, and model evaluation.

\IncMargin{1em}

\begin{algorithm}[!b]

	\SetKwData{new}{new}
	\SetKwFunction{embedsbert}{EmbedSBERT}
	\SetKwFunction{scalefeatures}{ScaleFeatures}
	\SetKwFunction{trainxgb}{TrainXGBoost}
	\SetKwFunction{evalmodel}{EvaluateModel}
	\SetKwFunction{sampleweights}{CalculateSampleWeights}
	
	\SetKwInOut{Input}{Input}
	\SetKwInOut{Output}{Output}
	
	\Input{ Dataset $D = \{(S_i, r_i, f_i)\}$ where $S_i$ are scripts, $r_i$ are ratings, $f_i$ are numerical features }
	\Output{ Trained XGBoost model $M$, best hyperparameters, performance metrics }
	\BlankLine
	
	$(X_{train}, X_{val}, X_{test}) \gets \text{split}(D, 0.70, 0.15, 0.15)$\; \label{split}\tcp*{70\% train / 15\% val / 15\% test}
	
	$X_{train}^{embed} \gets \text{EmbedSBERT}(X_{train}^{\mathrm{scripts}})$\; \label{embed_train}
	$X_{val}^{embed} \gets \text{EmbedSBERT}(X_{val}^{\mathrm{scripts}})$\; \label{embed_val}
	$X_{test}^{embed} \gets \text{EmbedSBERT}(X_{test}^{\mathrm{scripts}})$\; \label{embed_test}
	
	$X_{train}^{scaled} \gets \text{ScaleFeatures}(X_{train}^{\mathrm{feat}})$\;
	$X_{val}^{scaled} \gets \text{ScaleFeatures}(X_{val}^{\mathrm{feat}})$\;
	$X_{test}^{scaled} \gets \text{ScaleFeatures}(X_{test}^{\mathrm{feat}})$\;
	
	$X_{train}^{combined} \gets [X_{train}^{embed}, X_{train}^{scaled}]$\; \label{combine_train}
	$X_{val}^{combined} \gets [X_{val}^{embed}, X_{val}^{scaled}]$\;
	$X_{test}^{combined} \gets [X_{test}^{embed}, X_{test}^{scaled}]$\; \label{combine_test}
	
	$w \gets \text{CalculateSampleWeights}(X_{train}^{\mathrm{ratings}})$\; \tcp*{Class balancing weights}
	
	$M \gets \text{TrainXGBoost}(X_{train}^{combined}, X_{val}^{combined}, w)$\;
	
	$(RMSE_{val}, MAE_{val}, R^2_{val}) \gets \text{EvaluateModel}(M, X_{val}^{combined})$\;
	$(RMSE_{test}, MAE_{test}, R^2_{test}) \gets \text{EvaluateModel}(M, X_{test}^{combined})$\;
	
	\Return $(M, RMSE_{val}, MAE_{val}, R^2_{val}, RMSE_{test}, MAE_{test}, R^2_{test})$ \label{return_metrics}
	
	\caption{TrainXGBoostModel$(D)$}
	\label{alg:train_xgboost}
	
\end{algorithm}
\DecMargin{1em}


Key aspects of the training algorithm:

\begin{itemize}
    \item \textbf{Data Split} (Line~\ref{split}): 70\% training / 15\% validation / 15\% test using consistent random state (42)
    \item \textbf{SBERT Embeddings} (Lines~\ref{embed_train}-\ref{embed_test}): Generate 384-dimensional semantic representations for all three splits
    \item \textbf{Feature Scaling}: Standardize numerical features (movie length, year, decade, script statistics) using StandardScaler
    \item \textbf{Feature Combination} (Lines~\ref{combine_train}-\ref{combine_test}): Concatenate SBERT embeddings with numerical features to create 403-dimensional feature vectors
    \item \textbf{Class Balancing}: Compute sample weights inversely to class frequency to handle imbalanced rating distribution
    \item \textbf{XGBoost Training}: Train gradient boosting model with early stopping on validation set
    \item \textbf{Comprehensive Evaluation}: Report RMSE, MAE, and R² on both validation and test sets
\end{itemize}

\section{Class Balancing}

Since IMDb ratings follow a normal-like distribution with fewer movies at the extremes (excellent or terrible), we use sample weighting to ensure the model learns from all rating ranges adequately.

The weighting strategy is:
\begin{itemize}
    \item \textbf{Low (1-4)}: 5.25x weight (underrepresented)
    \item \textbf{Medium (4-6)}: 1.00x weight (baseline)
    \item \textbf{Good (6-8)}: 1.66x weight
    \item \textbf{Excellent (8-10)}: 13.24x weight (most underrepresented)
\end{itemize}

This ensures that rare rating ranges (especially excellent 8-10 movies) contribute proportionally to the loss function during training.

\section{Performance Metrics}

The trained XGBoost model achieved the following performance on the test set:

\begin{table}[h]
\centering
\begin{tabular}{lccc}
\toprule
\textbf{Dataset} & \textbf{RMSE} & \textbf{MAE} & \textbf{R²} \\
\midrule
Validation & 0.9419 & 0.7067 & 0.5659 \\
Test & 0.9881 & 0.7509 & 0.5762 \\
\bottomrule
\end{tabular}
\caption{Model Performance Metrics}
\label{tab:algo_performance}
\end{table}

\begin{itemize}
    \item \textbf{RMSE = 0.99}: Root Mean Square Error - typical prediction deviation is $\sim$1 point on 10-point scale
    \item \textbf{MAE = 0.75}: Mean Absolute Error - average absolute prediction error is 0.75 points
    \item \textbf{R² = 0.5762}: The model explains 57.62\% of the variance in IMDb ratings
\end{itemize}

\section{Complete Source Code}

Source code for the algorithms presented in this chapter can be found in Appendix B. The complete implementation includes:

\begin{itemize}
    \item SBERT embedding generation with chunking (\texttt{Codes/sbert\_embedding.py})
    \item XGBoost training pipeline (\texttt{Codes/xgboost\_trainer.py})
    \item Full model training script (\texttt{trainer.py})
    \item Prediction module (\texttt{predictor.py})
\end{itemize}

\endinput
